\section{Related Work}
\label{sec:related-work}

\subsection{Sandbox-Based Dynamic Analysis}

Automated malware sandboxes trace back to Cuckoo Sandbox~\cite{cuckoo},
whose architecture (a controller orchestrating detonation inside a
disposable virtual machine, with an in-guest agent reporting API calls
and artifacts back to the host) remains the standard pattern.
CAPE~\cite{cape} (Config And Payload Extraction) forked this lineage
specifically to add automated unpacking and configuration extraction
for known malware families, and ships with a large community-maintained
library of behavioral signatures, each optionally tagged with MITRE
ATT\&CK technique identifiers. This pipeline builds directly on CAPE
for dynamic analysis (\S\ref{sec:dynamic}), and treats its
community-maintained signature-to-technique mappings as a starting
hypothesis to verify rather than ground truth to trust outright, a
distinction substantiated concretely in \S\ref{sec:evaluation}: several
of CAPE's own signature mappings were checked against their raw
evidence and found to be systematically over-broad for the samples
they fired on, an inspection step this pipeline performs and most
sandbox deployments skip.

\subsection{Static ATT\&CK Mapping}

Mapping static artifacts (imports, strings, YARA matches) directly to
ATT\&CK techniques is common practice in commercial EDR products and in
community tooling. The closest open-source counterpart is
Mandiant's capa~\cite{capa}, which matches curated rules of imports,
strings, and byte patterns against a binary to report the capabilities
it implements, many of them pre-mapped to ATT\&CK. Most capa rules,
like most EDR indicator lookups, are single-indicator matches: one
import name or string pattern is enough to fire a rule. This project
deliberately diverges from that pattern in its confidence model
(\S\ref{sec:static-rules}): a single generic indicator (an import used
for many legitimate purposes) is treated as weaker evidence than a
\emph{coherent combination} of indicators that only makes sense
together for the claimed technique. Concretely, \technique{OpenProcess}
alone is not evidence of process injection (\technique{T1055}), but
\technique{OpenProcess} together with \technique{VirtualAllocEx} and
\technique{WriteProcessMemory} is, because the three APIs compose into
one recognizable primitive with no other common purpose. This trades
some of capa's breadth (capa ships hundreds of rules covering many
platforms and capability classes) for narrower, evidence-gated
precision on the Windows techniques this pipeline does cover, a
trade-off made explicit and measured directly in
\S\ref{sec:fp-audit}.

\subsection{Automated Extraction from Threat Reports}

A parallel line of work extracts ATT\&CK techniques not from a binary
but from unstructured analyst prose: TTPDrill~\cite{ttpdrill} parses
threat-intelligence reports with a custom ontology and NLP pipeline to
recover attacker actions; MITRE Engenuity's TRAM~\cite{tram} and
rcATT~\cite{rcatt} both train text classifiers to tag CTI report
paragraphs with likely ATT\&CK technique labels. These tools solve a
genuinely different problem: they help an analyst who already has a
\emph{written report} turn it into structured, machine-readable
technique labels faster. This pipeline instead derives its technique
labels directly from a sample's own binary and runtime behavior, with
no report as an intermediate step, which is what lets every finding
carry a specific piece of verifiable evidence (an import combination,
a signature's raw data field) rather than an NLP classifier's
confidence score over free text. The two approaches are complementary
rather than competing: a TRAM-style classifier could in principle
consume this pipeline's own written case studies (\S\ref{sec:case-studies})
as its input text, closing the loop from binary evidence to report to
re-extracted technique labels.

\subsection{Reputation Aggregators and Multi-Engine Scanners}

VirusTotal~\cite{virustotal} is the tool most incident responders reach
for first, and is a useful point of comparison because it operates at
the opposite end of the breadth-versus-depth trade-off this pipeline
makes. VirusTotal aggregates verdicts from dozens of antivirus engines
and sandbox integrations across a submission corpus of a scale no
individual research pipeline can match, which makes it excellent for
fast reputation triage: is this hash already known, and do other
vendors already flag it. What VirusTotal does not do is reconcile
those signals against each other for a single claim: each engine's
verdict and each sandbox's behavioral tags are surfaced side by side,
not cross-checked, so two contradictory technique tags from two
different integrations are both simply displayed, with no mechanism
resolving which one the raw evidence actually supports. This pipeline
targets that specific gap rather than VirusTotal's breadth: far fewer
samples, but every technique claim traced to the exact import
combination, signature data field, or string pattern that produced it,
cross-source and cross-level agreement modeled explicitly
(\S\ref{sec:reconciliation}), and a multi-stage sample's dropped
components analyzed and attributed individually rather than folded
into one aggregate verdict for the parent hash
(\S\ref{sec:resubmission}). The two tools are not substitutes: a
VirusTotal lookup is the right first step for triage at scale, and this
pipeline is the right next step once a specific sample needs a
defensible, evidence-linked technique mapping rather than a reputation
score.

\subsection{Threat Intelligence Interchange Standards}

STIX~2.1~\cite{stix21} (Structured Threat Information Expression) and
MISP~\cite{misp} (Malware Information Sharing Platform) are the two
interchange formats this pipeline exports to. STIX defines a graph of
typed objects (indicators, malware, attack patterns, relationships)
intended for machine-to-machine sharing; MISP is both a data model and
a running platform most organizations already operate. Exporting to
both, rather than only producing a human-readable report, is what
makes pipeline output directly consumable by a security operations
team's existing tooling instead of requiring manual transcription.

\subsection{Positioning}

This thesis does not propose a novel detection primitive: every
individual signal used (an API import, a string pattern, a sandbox
behavioral signature) is a well-established technique in malware
analysis practice, and tools like capa and CAPE already surface many
of the same raw signals. The contribution is in the
\emph{integration}: a confidence reconciliation model that is
verifiably more than the sum of its sources (\S\ref{sec:reconciliation}),
a mechanism for correctly attributing multi-stage malware's distributed
capability (\S\ref{sec:resubmission}), and, argued as the most
defensible contribution of the three, a validation discipline in which
every claim the pipeline makes is checked against real evidence before
being trusted, documented with enough specificity that the check itself
is reproducible (\S\ref{sec:evaluation}). Positioned against the two
poles surveyed above, this pipeline sits deliberately between a
broad, unreconciled aggregator like VirusTotal and a narrow,
single-sample manual reverse-engineering report: automated and
standards-exporting like the former, evidence-verified and
technique-specific like the latter.
